% a04_Bh_bulk1b  Proof of app:a04:lem:B2, part 1b: alias sums with the weight, cell averages, rows (I)+(II)

\subsubsection*{B2.3 The alias sums with the Dirichlet weight}

We use the Fourier representation from the proof of \cref{app:a04:lem:B0}: for $x,x'\in\Lambda_\tau$,
\[
\Gamma_\tau(x,x')=\int_{\mathbb T}\overline{\hat v^\tau_x(P)}\,\hat v^\tau_{x'}(P)\,\Mh_\tau(P)\dP,\qquad \Mh_\tau\defeq(b^{-1}+W_\tau)^{-1},\quad a_\tau(n)\defeq s_\tau(p)\hat C_\tau(p),
\]
with $p=P+2\pi n$, $n\in N_\tau$, and (E1)--(E4) as there; for the weight $f_D$, $s^f_\tau=s_{5\tau}(p)\varphi(\tau p)$, $a^f_\tau$, $\hat v^{f,\tau}_x$, $W^f_\tau$, $\Mh^f_\tau$ are the same expressions with $s^f_\tau$ in place of $s_\tau$ (the Poisson-summation derivation uses only that $f_\tau$ is $1$-periodic, so $Q_fC_\tau Q_f^*$ commutes with $\Z^3$-translations; $Q_fQ_f^*\ne1$ is not used). The continuum object is the same $\Gamma_0$ for both weights. Consequences recorded in that proof: $|\hat v^\tau_x|\le C(P^2+m'^2)^{-1}$, $\Mh_\tau\le\min(b,(\pi/2)^6(P^2+m'^2))$, $|\hat v^\tau_x-\hat v^0_x|,|W_\tau-W_0|\le C\tau(P^2+m'^2)^{-1/2}$, hence $|\Mh_\tau-\Mh_0|\le\Mh_\tau\Mh_0|W_\tau-W_0|\le C\tau(P^2+m'^2)^{3/2}$; for $n\ne0$, $c|n|\le|p|\le C|n|$; all uniform in $m'\in(0,m_0]$.

\begin{lemma}[alias sums]\label{app:a04:lem:Af}
For $\tau\le1$, $P\in\mathbb T$ and $\kappa\in[0,2)$, uniform weight:
\begin{equation}\label{app:a04:eq:star}
\textstyle\sum_{n\in N_\tau}|p|^\kappa|a_\tau(n)|\le C(P^2+m'^2)^{(\kappa-2)/2}+CS_\kappa\le C(P^2+m'^2)^{(\kappa-2)/2}.
\end{equation}
For the weight $f_D$ and $\tau=5^{-j}\le\frac15$, write $p=q+2\pi\ell/(5\tau)$ with $q\in\mathrm{BZ}_{5\tau}$, $\ell\in\{-2,\dots,2\}^3$, $q=P+2\pi n'$, $n'\in N_{5\tau}$ (a bijection $n\leftrightarrow(n',\ell)$); $\ell=0$ are the uniform aliases, $\ell\ne0$ the D-aliases ($|p|_\infty\ge\pi/(5\tau)$, so $\hat C_\tau(p)\le25\tau^2/4$, and $|p|\le\sqrt3\pi/\tau$). Then
\begin{enumerate}[label=\textup{(A\arabic*)},leftmargin=*]
\item $\sum_{\ell\ne0}|p|^\kappa|a^f_\tau(n)|\le C\tau^{2-\kappa}(1+\ln\tau^{-1})^3$;
\item $\sum_{n\in N_\tau}|p|^\kappa|a^f_\tau(n)|\le C(P^2+m'^2)^{(\kappa-2)/2}+CS_\kappa+C\tau^{2-\kappa}(1+\ln\tau^{-1})^3\le C(P^2+m'^2)^{(\kappa-2)/2}$;
\item $W^f_\tau(P)\ge(2/\pi)^6c_S^6(P^2+m'^2)^{-1}$, hence $\Mh^f_\tau\le\min(b,C(P^2+m'^2))$; $|\hat v^{f,\tau}_x|\le C(P^2+m'^2)^{-1}$; $W^f_\tau\le C(P^2+m'^2)^{-1}$;
\item $|\hat v^{f,\tau}_x-\hat v^0_x|+|W^f_\tau-W_0|\le C\tau(P^2+m'^2)^{-1/2}$; $|\Mh^f_\tau-\Mh_0|\le C\tau(P^2+m'^2)^{3/2}$; $|\nabla^{(\tau)}\hat v^{f,\tau}_x-\partial\hat v^0_x|\le C\tau(1+\ln\tau^{-1})^3$, where $\nabla^{(\tau)}$ multiplies the $n$-th term by $d_\tau(p_\mu)\defeq(e^{-i\tau p_\mu}-1)/\tau$ and $\partial$ by $-ip_\mu$.
\end{enumerate}
\end{lemma}
\begin{proof}
\eqref{app:a04:eq:star}: the $n=0$ term is $\le|P|^\kappa\frac{\pi^2}4(P^2+m'^2)^{-1}$ by (E1)--(E2), and the aliases give $\sum_{n\ne0}\rho(n)C|n|^{\kappa}|n|^{-2}=CS_\kappa$ (E3). (A1): each D-alias term is $\le(\sqrt3\pi/\tau)^\kappa(25\tau^2/4)\sigma(q)$ by (E1) at spacing $5\tau$ ($|s^f_\tau(p)|\le|s_{5\tau}(q)|\le\sigma(q)$, as in the Dirichlet-weight paragraph of \cref{app:a04:lem:B0}), and $\sum_{\ell\ne0}\sum_{n'\in N_{5\tau}}\sigma(q)\le124\big[1+\sum_{\emptyset\ne A\subset\{1,2,3\}}\prod_{\mu\in A}\sum_{1\le|n'_\mu|\le1/(10\tau)}\frac\pi{2|n'_\mu|}\big]\le C(1+\ln\tau^{-1})^3$. (A2): the $\ell=0$ terms are bounded as in \eqref{app:a04:eq:star} ($|s^f|\le\sigma(p)\le\rho(n)$); add (A1), and $\tau^{2-\kappa}(1+\ln\tau^{-1})^3\le C_\kappa$. (A3): the $n=0$ term of $W^f$ with $|s^f_\tau(P)|\ge(2/\pi)^3c_S^3$ (valid for $|\tau P_\mu|\le\pi/5$, i.e.\ $\tau\le\frac15$) and (E2); the upper bounds from (A2) at $\kappa=0$. (A4): telescope $a^f_\tau(n)-a_0(n)=(s^f_\tau-s_0)\hat C_\tau+s_0(\hat C_\tau-\hat C_0)$. Uniform aliases: $|s^f_\tau(p)-s_0(p)|\le|s_{5\tau}(p)-s_0(p)||\varphi|+|s_0(p)||\varphi(\tau p)-1|\le C\min(1,\tau|p|)\sigma(p)$ ((E4) at spacing $5\tau$ and $|\varphi_1(\vartheta)-1|\le C|\vartheta|$), and (E4) for $\hat C$, give $|a^f_\tau(n)-a_0(n)|\le C\tau\sigma(p)(|p|^2+m'^2)^{-1/2}$, summed $\le C\tau(P^2+m'^2)^{-1/2}+C\tau S_1$, as for the uniform weight. D-aliases: $\sum|a^f_\tau(n)|\le C\tau^2(1+\ln\tau^{-1})^3$ by (A1) at $\kappa=0$, and their continuum partners $a_0(n)$ with $|n|\ge1/(10\tau)$: $\sum_{|n|\ge N}\rho(n)|n|^{-2}\le CN^{-2}(1+\ln N)^2$. Continuum-only aliases $n\notin N_\tau$ ($|n|\ge(3\tau)^{-1}$): $\le C\tau$ as in \cref{app:a04:lem:B0}(ii). All are $\le C\tau\le C\tau(P^2+m'^2)^{-1/2}$ on $\mathbb T$. For $W^f-W_0$ there is one more factor $|s|\le1$; $|\Mh^f_\tau-\Mh_0|\le\Mh^f_\tau\Mh_0|W^f_\tau-W_0|$. Derivative clause: $|d_\tau(\vartheta)|\le|\vartheta|$, $|d_\tau(\vartheta)+i\vartheta|\le\tau\vartheta^2/2$; so $|\nabla^{(\tau)}\hat v^{f,\tau}_x-\partial\hat v^0_x|\le\sum_{n\in N_\tau}[|a^f_\tau-a_0||p|+|a_0|\tau p^2/2]+\sum_{n\notin N_\tau}|a_0||p|$; the uniform aliases give $C\tau\sum_{n\in N_\tau}\sigma(p)[|p|(|p|^2+m'^2)^{-1/2}+p^2(|p|^2+m'^2)^{-1}]\le C\tau(1+\sum_{N_\tau\setminus0}\rho(n))\le C\tau(1+\ln\tau^{-1})^3$, the D-aliases $C\tau(1+\ln\tau^{-1})^3$ by (A1) at $\kappa=1$, and $\sum_{|n|\ge N}\rho(n)|n|^{-1}\le CN^{-1}(1+\ln N)^2$ the rest.
\end{proof}

\subsubsection*{B2.4 Rows $(\mathrm I)$ and $(\mathrm{II})$: the translation-invariant part with rate}

\begin{lemma}[cell averages]\label{app:a04:lem:cell}
Let $\langle F\rangle_\cell\defeq\tau^3\sum_{x\in\Lambda_\tau\cap[0,1)^3}F(x,x)$. With $Z_\tau(P)\defeq\sum_{n\in N_\tau}|s_\tau|^2\hat C_\tau^3\Pih_\tau(p)$, $W_{2,\tau}\defeq\sum_n|s_\tau|^2\hat C_\tau^2$, $Y_\tau\defeq\sum_n|s_\tau|^2\hat C_\tau^2\Pih_\tau(p)$,
\[
\langle(\mathrm I)\rangle_\cell=2\int_{\mathbb T}\Mh_\tau Z_\tau\dP,\qquad\langle(\mathrm{II})\rangle_\cell=-\int_{\mathbb T}\Mh_\tau^2W_{2,\tau}Y_\tau\dP,
\]
and the same with $s^f_\tau$, $\Mh^f_\tau$ for the weight $f_D$. The continuum ($\tau=0$) expressions are the same with $N_0=\Z^3$, $s_0$, $\hat C_0$, $\Pih_0$.
\end{lemma}
\begin{proof}
$(\Gamma\Pi_CC)(x,x)=\sum_{y,z}\tau^6\Gamma(x,y)\Pi_C(y-z)C(z-x)=\int_{\mathbb T}\Mh_\tau\overline{\hat v_x(P)}\sum_na_\tau(n)\hat C_\tau(p)\Pih_\tau(p)e^{-ip\cdot x}\dP$ (sum the translation-invariant lines in Fourier). Averaging over $x\in\Lambda_\tau\cap[0,1)^3$ with weight $\tau^3$ kills the off-diagonal aliases, $\tau^3\sum_xe^{i(p-p')\cdot x}=\tau^3\sum_xe^{2\pi i(n-n')\cdot x}=\delta_{nn'}$ ($n,n'\in N_\tau$). For $(\mathrm{II})$: $(\Gamma\Pi_C\Gamma)(x,x)=\iint\Mh(P)\Mh(P')\overline{\hat v_x(P)}[\sum_{y,z}\tau^6\hat v_y(P)\Pi_C(y-z)\overline{\hat v_z(P')}]\hat v_x(P')$; the bracket is $\sum_{n,n'}a_\tau(n)\overline{a_\tau(n')}\Pih_\tau(p)\tau^3\sum_ze^{-i(p-p')\cdot z}$, and $\sum_{z\in\Lambda_\tau}\tau^3e^{-i(p-p')\cdot z}$ with $p,p'\in\mathrm{BZ}_\tau$ forces $P=P'$ (as a $\delta$ on $\mathbb T$) and $n=n'$; the cell average of $|\hat v_x|^2$ then gives $W_{2,\tau}$. The weight enters $s$ only, and the Poisson step is weight-free.
\end{proof}

\begin{theorem}[rows $(\mathrm I)+(\mathrm{II})$]\label{app:a04:thm:R}
For both weights, uniformly in $b\in[\mathfrak a(1-5^{-2}),\mathfrak a]$ and $m',\mu'\in(0,m_0]$ (for $f_D$: $\tau=5^{-j}\le\frac15$),
\[
\big|\langle(\mathrm I)+(\mathrm{II})\rangle_{\cell,\tau}-\langle(\mathrm I)+(\mathrm{II})\rangle_{\cell,0}\big|\le C\tau(1+m'^{-1}),\qquad
\big|\langle(\mathrm I)+(\mathrm{II})\rangle_{\cell,0}(b)-\langle(\mathrm I)+(\mathrm{II})\rangle_{\cell,0}(b')\big|\le C|b-b'|.
\]
Hence the diagonal increment obeys $\big|[(\mathrm I)+(\mathrm{II})]_{t/5}(\mathfrak a_{j+1})-[(\mathrm I)+(\mathrm{II})]_t(\mathfrak a_j)\big|\le Ct(1+m'^{-1})+Ct^2\le C't(1+m'^{-1})$.
\end{theorem}
\begin{proof}
\emph{(i) Majorants.} $n=0$ ($p=P\in\mathbb T$): $|s|\le1$, $\hat C\le C(P^2+m'^2)^{-1}$, $|\Pih|\le C(|P|+m_0)$ ($\Pi$-a), $\Mh\le C(P^2+m'^2)$ ((A3) for $f_D$). So $|\Mh Z^{(0)}|\le C(|P|+m_0)(P^2+m'^2)^{-2}$ and $\int_{\mathbb T}\le C\int_0^{\sqrt3\pi}(\varrho+m_0)\varrho^2(\varrho^2+m'^2)^{-2}d\varrho\le C(1+\ln(1/m')+m_0/m')\le C(1+m'^{-1})$. Uniform aliases $n\ne0$: $|s|^2\le\rho(n)^2$, $\hat C\le\frac34|n|^{-2}$, $|\Pih|\le C|n|$, so $|\Mh Z^{(n)}|\le C\rho(n)^2|n|^{-5}$ and $\sum_{n\ne0}\le CS_0$. D-aliases ($f_D$): $|\Mh^fZ^{f,(n)}|\le C\sigma(q)^2(25\tau^2/4)^3C|p|\le C\tau^5\sigma(q)^2$, summed $\le C\tau^5(1+\ln\tau^{-1})^3$ by the count in (A1); likewise $W^f_2$ ($\le C\tau^4(1+\ln\tau^{-1})^3$) and $Y^f$ ($\le C\tau^3(1+\ln\tau^{-1})^3$).
\emph{(ii) Differences at a common $n\in N_\tau$.} Telescope the product $\Mh|s|^2\hat C^3\Pih$ one factor at a time; each factor's difference is $\le C\tau(1+|p|)$ times its majorant: $|s_\tau-s_0|\le C\tau|p|\sigma$ (E4; for $f_D$ and $\ell=0$ the bound of the proof of (A4)), $|\hat C_\tau-\hat C_0|\le C\tau(|p|^2+m'^2)^{-1/2}=C\tau(|p|^2+m'^2)^{1/2}(|p|^2+m'^2)^{-1}$ (E4), ($\Pi$-b), and $|\Mh_\tau-\Mh_0|\le C\tau(P^2+m'^2)^{3/2}=C\tau(P^2+m'^2)^{1/2}(P^2+m'^2)$. For $n=0$ this gives $\int_{\mathbb T}C\tau(1+|P|)(|P|+m_0)(P^2+m'^2)^{-2}\le C\tau(1+m'^{-1})$; for $n\ne0$, $\sum C\tau(1+|n|)\rho(n)^2|n|^{-5}\le C\tau$. D-aliases: bound the $\tau$-term by its majorant of (i) ($O(\tau^3(1+\ln\tau^{-1})^3)$ at worst) and its continuum partner by the tail $|n|\ge1/(10\tau)$: $|\Mh_0Z_0^{(n)}|\le C\rho(n)|n|^{-5}$, $\sum_{|n|\ge1/(10\tau)}\le C\tau^2$.
\emph{(iii) Continuum-only aliases} $n\notin N_\tau$ ($|n|\ge(3\tau)^{-1}$): $|\Mh_0Z_0^{(n)}|\le C\rho(n)|n|^{-5}$, $\sum_{|n|\ge1/(3\tau)}\le C\tau^2$.
\emph{(iv) $(\mathrm{II})$.} The same bookkeeping for $\Mh^2W_2Y$ with the majorants $\Mh^2\le C(P^2+m'^2)^2$, $W_2^{(0)}\le C(P^2+m'^2)^{-2}$, $Y^{(0)}\le C(|P|+m_0)(P^2+m'^2)^{-2}$, alias parts $\le CS_0$ (and the D-alias bounds of (i)); every product of majorants is $\le C[(|P|+m_0)(P^2+m'^2)^{-2}+1]$, integrable with $\int\le C(1+m'^{-1})$, and each factor's difference again costs $C\tau(1+|p|)$.
\emph{(v) $b$-dependence at $\tau=0$.} $\partial_b\Mh=b^{-2}\Mh^2\le C(P^2+m'^2)^2$, so $|\partial_b\langle(\mathrm I)\rangle_0|\le C\int(P^2+m'^2)^2(|P|+m_0)(P^2+m'^2)^{-3}\le C$ and $|\partial_b\langle(\mathrm{II})\rangle_0|\le C\int(P^2+m'^2)^3(|P|+m_0)(P^2+m'^2)^{-4}\le C$ (no infrared loss). The continuum expressions are weight-free (same $\Gamma_0$).
\emph{(vi) Diagonal increment.} Insert $\langle\cdot\rangle_0(\mathfrak a_{j+1})$ and $\langle\cdot\rangle_0(\mathfrak a_j)$; $|\mathfrak a_j-\mathfrak a_{j+1}|\le Ct^2$. Each member is compared with the continuum with its own weight and spacing ($t$ or $t/5$).
\end{proof}
